%!TEX root = ./main.tex

\section[Wrap-up]{Wrap-up}

% SECTION TIMING: 45:00--52:00 (3 frames). Summary, then the honest state of the field,
% then collect the exit ticket. Do not start a new idea here.

% Rewritten 2026-09-13: was a per-paper comparison table; now the three ideas of the
% lecture in the order they were taught.
\begin{frame}{Three things to take home}
  \vspace{0.6em}
  \begin{enumerate}\itemsep14pt
    \item \textbf{One waveform, two jobs.} The signal a base station already sends
      comes back as an echo, and that echo is measurement nobody had to upload.
    \item \textbf{The two jobs want opposite things.} Sensing wants a wide, steady,
      predictable signal; data wants an efficient, bursty, random one. So there is a
      curve of compromises, not a fix.
    \item \textbf{On OFDM, picking a point on that curve is a scheduling decision.}
      Which cells of the grid carry bits, which carry echoes, and who picks first.
  \end{enumerate}

  \vspace{0.9em}
  \takeaway{ISAC is not a new radio.\\It is a new question asked of the radio we already have: who gets each cell?}
\end{frame}

\begin{frame}{What is actually unfinished}
  \vspace{0.5em}
  \begin{columns}[T,onlytextwidth]
    \column{0.47\textwidth}
      \large\textbf{\textcolor{softgreen}{Settled enough to build}}

      \vspace{0.5em}
      \small
      \begin{itemize}\itemsep4pt
        \item The tradeoff is understood.
        \item OFDM-based ISAC works.
        \item The RE grid can be allocated.
      \end{itemize}
    \column{0.47\textwidth}
      \large\textbf{\textcolor{coreorange}{Still open}}

      \vspace{0.5em}
      \small
      \begin{itemize}\itemsep4pt
        \item A killer application.
        \item Security of a shared waveform.
        \item AI for waveform design.
        \item Space--air--ground reach.
      \end{itemize}
  \end{columns}

  \vspace{0.9em}
  \qa{Which of these is not an engineering problem?}
     {The first. A waveform can be optimised; a market cannot. ISAC has a
      technology looking for the application that pays for it.}

  \glossline{SAGIN = Space--Air--Ground Integrated Network}
\end{frame}

% TEACHING NOTE (2:30): collect on paper or in the LMS. Require a marked point and one
% sentence. Do not accept "both" as an answer to part 3.
\begin{frame}{Before you go: pick a point on the curve}
  \vspace{0.2em}
  {\small Two minutes, on paper, handed in on the way out. Same curve as before.\par}
  \begin{center}
  \begin{tikzpicture}[scale=0.85,transform shape]
    \draw[-{Latex[length=2.2mm]},softgray] (0,0) -- (5.6,0)
      node[anchor=north east,font=\scriptsize] {sensing error $\rightarrow$ worse};
    \draw[-{Latex[length=2.2mm]},softgray] (0,0) -- (0,3.0)
      node[anchor=south,font=\scriptsize] {rate};
    \draw[very thick,draw=ranblue] (0.6,2.7) .. controls (2.0,2.5) and (2.8,1.6) .. (5.0,1.1);
  \end{tikzpicture}
  \end{center}

  \vspace{0.6em}
  \begin{enumerate}\itemsep6pt
    \item Mark where a \textcolor{ranblue}{self-driving car} should sit. Mark where a
      \textcolor{coreorange}{home Wi-Fi link} should sit.
    \item Name one resource the two designs would allocate differently.
    \item Answer in one sentence:\\
      \textbf{what does sensing cost the bits in your example?}
  \end{enumerate}
\end{frame}
